\chapter{System Design}

\section{System Architecture}
The application follows a state-based architecture managed by a Finite State Machine (FSM). This ensures a linear and controlled training flow. The main states are:
\begin{enumerate}
    \item \textbf{Welcome State:} Handles environment loading and user orientation.
    \item \textbf{Video State:} Manages the custom OpenCvSharp video decoder to render frames to a texture in real-time.
    \item \textbf{Quiz State:} Manages the logic for the 12-question assessment, handling user input and scoring.
    \item \textbf{Finished State:} Aggregates results and provides a feedback loop.
\end{enumerate}

\section{Activity Diagram}
The following diagram illustrates the user flow through the application:

\begin{figure}[h]
    \centering
    % Placeholder for Activity Diagram
    \fbox{\begin{minipage}{10cm}
        \vspace{5cm} \centering [Insert Activity Diagram Here] \vspace{5cm}
    \end{minipage}}
    \caption{Activity Diagram of the Safety Training System}
    \label{fig:activity_diagram}
\end{figure}

\section{Class Diagram}
The core logic is encapsulated in the \texttt{Program} class, which manages the main loop and sub-systems.
\begin{itemize}
    \item \texttt{Program}: Main entry point, FSM manager.
    \item \texttt{SeatCollider}: Struct for collision detection logic.
    \item \texttt{QuizQuestion}: Data class for holding question text, options, and correct answers.
\end{itemize}

\section{Data Flow}
User input (Controller/Keyboard) $\rightarrow$ Input Handler $\rightarrow$ State Machine $\rightarrow$ Rendering Engine $\rightarrow$ Display.
